\documentclass[10pt,a4paper]{amsart}
\usepackage[margin=19mm]{geometry}
\usepackage{amsmath,amssymb,mathtools}
\usepackage[scr=rsfs,cal=euler]{mathalfa}
\usepackage{xfrac}
\usepackage[unicode,hidelinks,bookmarks=false]{hyperref}
\newcommand{\ie}{i.e.\ }
\newcommand{\cf}{cf.\ }
\newcommand{\resp}{respectively\ }
\newcommand{\Subsection}{\subsection}
\providecommand{\nobreakdash}{\nobreak\hbox{-}}
\setlength{\emergencystretch}{2em}
\multlinegap0pt
\usepackage{bm}
\usepackage{bbm}
\usepackage{dsfont}
\usepackage{xspace}
\usepackage{cancel}
\usepackage{mathtools}
\usepackage{cleveref}
\usepackage{amsthm}
\makeatletter
\@ifundefined{theorem}{\newtheorem{theorem}{Theorem}}{}
\@ifundefined{lemma}{\newtheorem{lemma}[theorem]{Lemma}}{}
\makeatother
\title[On the tightness of left-invariant contact structures]{On the tightness of left-invariant contact structures}
\author{Eugenio Bellini}
\date{Source: arXiv:2504.06085v3 (CC BY 4.0)}
\begin{document}
\maketitle

\noindent\emph{Source and licence.} On the tightness of left-invariant contact structures. Version arXiv:2504.06085v3 (CC BY 4.0), distributed under the Creative Commons Attribution 4.0 International licence, as recorded in the arXiv licence metadata for this exact version and shown on the abstract page https://arxiv.org/abs/2504.06085v3. Full rights notice: licenses/arxiv-2504.06085-CC-BY-4.0.txt.\par\smallskip

\noindent\textit{Editorial scope.} Counted unit: the Lemma whose source label is \texttt{lem:orthogonal\_groups} in the article (source0.tex lines 241--247, source label \texttt{lem:orthogonal\_groups}), delivered together with the article's own complete original proof of it (source0.tex lines 249--263, one \texttt{proof} environment of 15 lines). No mathematics is written, repaired or re-derived: statement and proof are the article's own text line for line. Required context retained uncounted: 4cor:groups (lines 217--219). Every definition, display and label of those lines is retained unchanged. Omitted from this package and not redistributed: 1--216, 220--240, 248--248, 264--536. Nothing inside a retained excerpt is omitted or altered: every retained body is delivered exactly as the article's lines are written, so no omission receipt of any kind accompanies this package. Citations: the retained excerpts carry no citation command at all, so no bibliography entry of the article is transcribed and no reference is redistributed. Reference adaptation: none was needed and none was made; every reference inside the delivered unit resolves inside it, so no unresolved reference or citation is produced. Printed slips kept literal: no printed word, symbol, hypothesis or number of the counted statement or of its proof was altered. Layout and class: the article's own class is \texttt{amsart} with packages including amsaddr, amsmath, amssymb, amsfonts, mathrsfs, bbold, babel, inputenc, fontenc, mathtools, lmodern, comment, cancel, geometry; the wrapper is the lane's pinned \texttt{amsart} editorial wrapper at 10pt on A4, which restates the theorem-like environments the retained text needs and adds the article's own macro definitions that the retained text needs (none), with extra wrapper packages (bm, bbm, dsfont, xspace, cancel, mathtools, cleveref). The delivered numbering is the unit's own. Assets: no retained span contains a figure environment, a graphics-inclusion command, a PDF-page inclusion, a table or a TikZ picture; no image file is copied into this package, the package declares no asset dependency and no image file is redistributed with it. Licence: version arXiv:2504.06085v3 of this article is distributed under the Creative Commons Attribution 4.0 International licence, as recorded in the arXiv licence metadata for this exact version and shown on the abstract page https://arxiv.org/abs/2504.06085v3. A full rights notice is delivered at licenses/arxiv-2504.06085-CC-BY-4.0.txt. Characters: every retained excerpt is pure ASCII, so the delivered engine, XeLaTeX with its default Latin Modern text font, reports no missing character.
\begin{lemma}\label{4cor:groups}
    Let $G$ be a Lie group endowed with a left-invariant Riemannian metric. Let $H\subset G$ be an immersed Lie subgroup. If the normal exponential map $\exp^{\perp}:TH^\perp\to G$ is an immersion then it is a covering map.
\end{lemma}

\begin{lemma}\label{lem:orthogonal_groups}
	Let \(G\) be a simply connected Lie group endowed with a left-invariant Riemannian metric. Assume the existence of an immersed codimension-one subgroup \(H\) and a totally geodesic subgroup \(H_3\), orthogonal to each other. Then the map
	\[
	p:H \times H_3 \longrightarrow G, \qquad p(h,h_3)=hh_3
	\]
	is a diffeomorphism.
\end{lemma}

\begin{proof}
	Let \(X\) be a left-invariant vector field spanning \(TH_3\). Since \(H_3\) is totally geodesic and orthogonal to \(H\), the vector field \(X\) is geodesic and orthogonal to \(H\). Therefore, the normal exponential map can be written as
	\[
	\exp^\perp: TH^{\perp} \simeq H \times \mathbb{R} \to G, \qquad \exp^\perp(h,t) = e^{tX}(h).
	\]
	Because \(X\) and \(H\) are transverse, \(\exp^\perp\) is an immersion. By \cref{4cor:groups}, it is also a covering map. Since \(G\) is simply connected, \(\exp^\perp\) is a diffeomorphism. It follows that the map
	\[
	\psi: \mathbb{R} \to H_3, \qquad \psi(t) =\exp^\perp(1_G,t)=e^{tX}(1_G),
	\]
	where $1_G$ denotes the identity element of $G$, is a diffeomorphism as well. Finally, the map \(p: H \times H_3 \to G\) can be expressed as 
	\[
	p(h,h_3)=hh_3=e^{\psi^{-1}(h_3)X}(h)=\exp^{\perp}(h,\psi^{-1}(h_3)),
	\]
	where the second equality follows from left-invariance of $X$ and definition of $\psi$. Since $\exp^{\perp}$ and $\psi$ are diffeomorphisms, we conclude the proof.
\end{proof}

\end{document}
